\section{Quadratic programming and complementarity}\label{sec:quadratic}

At a minimum of a convex quadratic over a polyhedron, each inequality
has either zero slack or zero multiplier. If we knew which entry to set
to zero in every pair, we could recover the point and its multipliers
by solving linear equations and inequalities. We find those zeros by
applying the pricing method of Section~\ref{sec:lp} to the more general
monotone linear complementarity problem~\cite{CD,CPS}.

Given \(M\in\RR^{N\times N}\) and \(q\in\RR^N\), the linear
complementarity problem asks for \(z,w\in\RR^N\) satisfying
\begin{equation}\label{eq:lcp}
 z\ge0,\qquad w=Mz+q\ge0,\qquad z^\top w=0 .
\end{equation}
Here \(M\) is \emph{monotone}, meaning \(M+M^\top\succeq0\).
Nonnegativity makes the last equality equivalent to \(z_iw_i=0\)
for every \(i\).

\begin{theorem}[Monotone linear complementarity]\label{thm:lcp}
There is a deterministic strongly polynomial algorithm
that, for any monotone linear complementarity
problem~\eqref{eq:lcp},
returns a solution or an infeasibility certificate \(u\ge0\)
satisfying \(M^\top u\le0\) and \(q^\top u<0\).
\end{theorem}

For convex quadratic programming the earlier bit-polynomial algorithm
is due to Kozlov, Tarasov and Khachiyan~\cite{KTK}. Interior-point
methods for monotone complementarity include those of Kojima, Mizuno
and Yoshise~\cite{KMY}. Granot and Skorin-Kapov obtained Tardos-type
bounds with matrix dependence for strictly convex quadratics~\cite{GSK},
and Cui, Li and Miao obtained such bounds for separable
quadratics~\cite{CLM}. V\'egh's strongly polynomial algorithms for
network flows with convex costs include separable quadratic costs~\cite{VQ}.

Monotonicity enters at two distinct steps. It ensures that feasibility
of the linear constraints in~\eqref{eq:lcp} suffices for existence of a
complementary solution. To \emph{find} one, we also use it to construct
a convex set that approximates every scaled, normalized solution
and has small complementarity gap throughout. The midpoint argument of
Lemma~\ref{lem:mid} then makes the coordinate labels unambiguous.

\subsection{Linear feasibility and existence of a solution}

By Cottle's theorem~\cite{Cottle}, a monotone LCP has a solution
whenever its linear constraints are feasible. For the proof, we first
record the attainment argument for convex quadratics, which will also
give the unboundedness certificate in Corollary~\ref{cor:qp}.

Consider \(\tfrac12 x^\top Gx+c^\top x\) on a nonempty polyhedron
\(\{x:Ax\le b\}\), where \(G\) is symmetric positive semidefinite.
The Frank--Wolfe attainment theorem~\cite{FW} in this setting follows
from the following alternative. There may be a direction

\begin{equation}\label{eq:baddir}
 d:\quad Ad\le0,\qquad Gd=0,\qquad c^\top d<0,
\end{equation}
giving rays of descent to \(-\infty\). Otherwise
\(c=Gv-A^\top\mu\) for some \(v,\mu\) with \(\mu\ge0\), by Farkas.
Writing \(G=T^\top T\), the objective is then
\[
 \tfrac12\|T(x+v)\|_2^2-\mu^\top Ax+\mathrm{const}.
\]
The pairs \((T(x+v),\mu^\top Ax)\) attainable with \(Ax\le b\)
form a polyhedron, with second component bounded above.
So in this image space the objective has compact sublevel sets
on the polyhedron. There is an optimum.


Now minimize
\(z^\top w\), \(w=Mz+q\), subject to \(z,w\ge0\).
If feasible there is an optimum by attainment, since the objective
is a convex quadratic and bounded below by zero.
Any such optimum solves the LCP. Indeed, in the variable \(z\)
it also minimizes the linearization there, so by LP duality
there are multipliers \(u,v\ge0\) with
\[
 w+M^\top z=v+M^\top u,\qquad w^\top u=z^\top v=0.
\]
Dotting the stationarity equation (with all terms on the left)
with \(z-u\) gives
\[
 z^\top w+u^\top v+(z-u)^\top M(z-u)=0.
\]
All three terms are nonnegative, proving the claim.

\subsection{Labeling for complementarity}

First use Theorem~\ref{thm:lp} to test the linear constraints
\(-z\le0,\ -Mz\le q\). In the infeasible case, the Farkas
multipliers on \(-Mz\le q\) form the desired certificate \(u\).
In the feasible case with \(q=0\), take \(z=w=0\). For \(q\ne0\)
collect the scaled affine spaces in the single linear subspace
\[
 W=\{(z,Mz+\alpha q):z\in\RR^N,\ \alpha\in\RR\}.
\]
View these vectors \(y=(z,w)\) as rows when used with bases
and scalings. Use positive slots for \(z\), negative for \(w\),
so \(Q(y)=z^\top w\). Use as row basis \(B_0\)
a basis of the directions with \(\alpha=0\), say from unit vectors
in the \(z\)-space, followed by \((0,q)\).


Choose a solution \(y\), with \(\alpha=1\), of minimal support.
It is a circuit of \(W\). Indeed, if the subspace supported on
\(\supp y\) had dimension at least two, the linear functional
\(\alpha\) would vanish on a nonzero vector there. Move from \(y\)
in one of the two directions along this vector until a positive
entry first reaches zero. The resulting vector still has
\(\alpha=1\), all its entries are nonnegative, and it retains at
least one zero in each pair. It would therefore be a solution
of smaller support.

As before, apply calibration (Theorem~\ref{thm:cal})
and the extremal price theorem (Theorem~\ref{thm:prices}),
restarting price search on a calibration restart.
Use the same parameters~\eqref{eq:labelpars}.
At a queried price vector let \(B_*=LB_0D\)
be the true basis corresponding to the answer \(B\)
as in Theorem~\ref{thm:cal}.
Our task is to give label \(h\) at \(i\)
if slot \((i,h)\) maximizes surplus for such a solution circuit
(with its calibrated demands).
After scaling the circuit by \(D\) and normalizing, the entry
at such a slot is at least \(1/3\).
As in LP we will test for large entries using only \(B\).
To avoid conflicting labels, we again want
a convex set of candidates on which \(Q\) is uniformly small,
to use the midpoint argument (Lemma~\ref{lem:mid}).
We develop the constraints first with the true basis.

\paragraph{A convex constraint for a scaled solution.}
With the reciprocal scaling \(D\), the complementarity gap
\(Q(y)=z^\top w\) is unchanged. On the subspace with \(\alpha=0\),
the gap is \(z^\top Mz\ge0\). The first \(N\) rows of \(B_*=LB_0D\)
span this subspace after scaling because \(L\) is lower triangular.
The last diagonal entry of \(L\) is positive, so a vector with
\(\alpha>0\) has a positive coefficient on the last row.

Write the row coefficients as \(\xi=(v^\top,\lambda)\) with
\(v\in\RR^N\). Gathering the terms that contain \(\lambda\) gives
\[
 Q(\xi B_*)=v^\top A_*v-\lambda g_*(\xi),
\]
where \(A_*\) is symmetric positive semidefinite and \(g_*\) is
linear. In the scaled, normalized representation of a solution,
\(\lambda>0\) and \(Q(\xi B_*)=0\). Hence
\(g_*(\xi)=v^\top A_*v/\lambda\ge0\). Every such solution belongs to
the convex set
\[
 \lambda\ge0,\qquad g_*(\xi)\ge0,\qquad
 v^\top A_*v\le \lambda g_*(\xi).
\]
Indeed, \(\sqrt{v^\top A_*v}\) is a seminorm and the geometric mean of
the two nonnegative quantities on the right is concave. This is the
rotated second-order cone constraint in the coordinates of \(B_*\).

\paragraph{A label test on the rational basis.}
With \(B\), use again the decomposition
\(Q(\xi B)=v^\top A v-\lambda g(\xi)\),
taking \(A\) symmetric as before. Recall \(H=2^{100N^2}\),
\(\epsilon=H^{-10}\), and set
\(
 \widetilde A=A+H^{-50}I.
\)
Use candidates \(y=\xi B\)
as in~\eqref{eq:box}
(again with \(\|\xi\|_\infty\le H\),
\(-\epsilon\le y_s\le 2\) for each entry),
with the \emph{additional} constraints
\begin{equation}\label{eq:rot}
 \lambda\ge0,\qquad
 g(\xi)+\epsilon\ge0,\qquad
 v^\top \widetilde A v\le(\lambda+\epsilon)(g(\xi)+\epsilon).
\end{equation}
These height shifts let us work with only absolute accuracy on the
basis entries,
even when a solution has \(\lambda\) tiny or \(g_*(\xi)=0\).


Indeed, \(B\) is within \(H^{-100}\) of \(B_*\),
which is bounded by \(2\) entrywise. Hence \(A,g\) are within
\(20NH^{-100}\) of \(A_*,g_*\) coefficientwise, by taking dot
products. Thus \(\widetilde A\) is positive semidefinite,
so the constraints are convex together as before.
Also for \(\|\xi\|_\infty\le H\),
\[
 |g(\xi)|\le H^2,\qquad
 |g(\xi)-g_*(\xi)|\le H^{-90},
 \qquad |v^\top(\widetilde A-A_*)v|\le H^{-45}.
\]
All candidates satisfy
\[
 Q(y)\le \epsilon(\lambda+g(\xi)+\epsilon)\le H^{-1}.
\]
And for any solution \(y_0\), write
\(y_0D/\|y_0D\|_\infty=\xi B_*\), with
\(\|\xi\|_\infty\le H\).
To see that \(y=\xi B\) is a candidate, observe that it approximates \(\xi B_*\) within \(\epsilon^2\)
entrywise, and~\eqref{eq:rot} holds since in particular
\[
 (\lambda+\epsilon)(g(\xi)+\epsilon)
 \ \ge\ \lambda g_*(\xi)+\epsilon^2/2
 \ \ge\ v^\top \widetilde A v.
\]


Assign \(+\) at \(i\) if the candidate set contains a point with
\(y_{i+}\ge1/4\), and assign \(-\) otherwise. Only short rational
data enter this test. To carry it out using the bit-polynomial convex
QP algorithm~\cite{KTK}, divide by the positive quantity
\(\lambda+\epsilon\). Introduce
\(s=1/(\lambda+\epsilon)\), \(\bar\xi=s\xi\) and \(\bar y=sy\).
Multiply every linear constraint, including the label threshold, by
\(s\), and impose
\[
 s\ge\frac{1}{H+\epsilon},\qquad
 \bar\lambda+\epsilon s=1,\qquad
 \bar v^\top\widetilde A\bar v\le g(\bar\xi)+\epsilon s,
 \qquad \bar\xi=(\bar v^\top,\bar\lambda).
\]
The last inequality is~\eqref{eq:rot} multiplied by \(s^2\).
Conversely, dividing the barred variables by \(s>0\) recovers the
original constraints. This perspective transformation~\cite[\S2.3.3]{Boyd}
leaves one convex quadratic inequality on a bounded polyhedron.
Indeed, \(\bar\lambda\ge0\) gives \(s\le1/\epsilon\), and the scaled
box constraints bound \(\bar\xi,\bar y\). Test linear feasibility
first. On a nonempty polyhedron, minimize the left side minus the right
side of the quadratic inequality and accept exactly when its minimum
is nonpositive.

\begin{proof}[Conclusion of the proof of Theorem~\ref{thm:lcp}]
As in LP, consider the final uninterrupted price search, with demands given by
calibration. A solution circuit, scaled and normalized at a query,
has entry at least \(1/3\) at each slot of maximal surplus on
its support. A candidate therefore reaches \(1/4\) at this slot.
No candidate can reach the same threshold at the opposite slot by convexity and
Lemma~\ref{lem:mid}. As a result the labels are consistent
with the solution circuit's demands.


Write \(x\) for the prices returned by the extremal price algorithm.
In each pair allow only one slot to be positive, namely at the
sign of \(x_i\) (arbitrary choice when zero).
Return to the original system~\eqref{eq:lcp}
and impose the indicated zeros.
Each demand on the solution circuit's support is positive and
met, so these zeros retain the solution circuit. Theorem~\ref{thm:lp}
recovers a solution of the resulting linear system.

The number and sizes of the label tests obey the same bounds as in
Section~\ref{sec:lp}. Their convex QP data have short binary
encodings. The two LP calls on the original data check initial
feasibility and recover the solution. All steps are deterministic.
\end{proof}

A single solution \((z_*,w_*)\) also gives a linear description of
the entire solution set~\cite{Man}. On the region \(z,w\ge0\),
\(w=Mz+q\), the identity
\[
 z^\top w=z^\top w_*+z_*^\top w+(z-z_*)^\top M(z-z_*)
\]
splits the gap into three nonnegative terms. Their sum is zero
precisely when
\[
 z^\top w_*+z_*^\top w=0,\qquad (M+M^\top)(z-z_*)=0 .
\]
The second equivalence uses positive semidefiniteness of \(M+M^\top\).
Consequently the solution set is a polyhedron. To find a solution of
maximum support, test which entries can be positive by linear
feasibility. Averaging the witnesses gives a solution positive on
all those entries at once.

\subsection{Quadratic optimization and conic feasibility}

The first-order optimality conditions for a convex quadratic are
an affine variational inequality~\cite{KT,KuhHist}. Their classical
complementarity formulation~\cite{CD,CPS} yields the following
algorithm.

\begin{corollary}[Convex quadratic programming and affine variational inequalities]\label{cor:qp}
Both problems below are solvable in deterministic strongly
polynomial time for arbitrary real data.
\begin{enumerate}
 \item \emph{Affine variational inequalities.} Given a polyhedron \(P=\{x\in\RR^n:Ax\le b\}\)
 and monotone \(G\in\RR^{n\times n}\), \(c\in\RR^n\),
 find \(x\in P\) such that
 \((Gx+c)^\top(x'-x)\ge0\) for all \(x'\in P\), or report
 that none exists.
 \item \emph{Convex quadratic programming.} For such \(P\) and symmetric \(G\succeq0\), minimize
 \(\tfrac12 x^\top Gx+c^\top x\) over \(P\), returning an
 optimum and multipliers, or certifying infeasibility or
 unboundedness.
\end{enumerate}
\end{corollary}


\begin{proof}
By LP duality the variational inequality at \(x\in P\) holds
if and only if there are multipliers satisfying
\[
 Gx+c+A^\top\mu=0,\qquad
 \mu\ge0,\qquad
 \mu^\top(b-Ax)=0.
\]
Write \(x=Ru\) for \(R=[I\ {-I}]\), with \(u\ge0\).
The requirements on \(u,\mu\) are exactly~\eqref{eq:lcp}
for
\[
 z=\binom{u}{\mu},\qquad
 M=\begin{pmatrix}
  R^\top GR & R^\top A^\top\\ -A R & 0
 \end{pmatrix},\qquad
 q=\binom{R^\top c}{b}.
\]
Indeed, applying \(R^\top\) includes each entry of
\(Gx+c+A^\top\mu\) and its opposite in the first block
of \(w\), forcing zero. Monotonicity is
immediate, so apply Theorem~\ref{thm:lcp}.

For the quadratic objective, first test feasibility of \(P\) by
Theorem~\ref{thm:lp}. If \(P\) is empty this gives a Farkas
certificate. Otherwise retain a feasible point \(x_0\) and test the
linear conditions~\eqref{eq:baddir} on \(d\). Any such \(d\) gives
the ray \(x_0+td\), \(t\ge0\), along which the objective decreases
to \(-\infty\). If no such direction exists, the attainment argument
gives a minimum. Convexity makes the variational inequality necessary
and sufficient for optimality, so the complementarity solution yields
an optimizer and its multipliers. These reductions use a constant
number of layers of scalar computations and calls as in
Section~\ref{sec:arithmetic}.
\end{proof}

The LCP and convex QP algorithms and their certificates transfer to
ordered fields by the argument in Section~\ref{sec:lp}. Diagonalizing
the symmetric form by congruence shows that its positive
semidefiniteness persists in a real closure.

\paragraph{Projection.}
For any \(v\) and any nonempty polyhedron \(P\), minimizing
\(\|x-v\|_2^2\) over \(P\) gives its Euclidean projection.
If a polytope is given by vertices, barycentric variables express
membership by linear constraints. This includes the minimum-norm
point of a vertex-given simplex.
See Hong and Verma~\cite{HV} for earlier connections of
strong polynomiality questions for convex QPs and projection, and
De Loera, Haddock and Rademacher~\cite{DLHR}
for an LP reduction \emph{to} the minimum-norm problem in an
integer-input model, as well as exponential behavior of Wolfe's
combinatorial algorithm \cite{W76}.
Fujishige, Kitahara and V\'egh~\cite{FKV} give bounds using circuit
imbalance for update-and-stabilize methods for box-constrained least
squares.

\paragraph{A conic consequence.}
Another short application, using perspective again,
is to a special case of second-order cone programming
(see~\cite{Lobo,Boyd} for background and applications).

\begin{corollary}[Lorentz constraints with proportional quadratic parts]\label{cor:lor}
One can deterministically decide, in strongly polynomial time,
real feasibility of linear weak inequalities together with a single
Lorentz constraint
\begin{equation}\label{eq:lor}
 \|F(x)\|_2\le h(x)
\end{equation}
for real affine maps \(F,h\) with vector and scalar output, respectively,
and return a point in case of feasibility.

More generally we can decide feasibility, within the same bound,
for arbitrarily many Lorentz constraints if the Hessians
\(\nabla^2(\|F(x)\|^2_2-h(x)^2)\)
(one for each constraint)
span at most one dimension
\textup{(}e.g., ball constraints of the form
\(\|x-c\|_2\le r\) for given \(c,r\)\textup{)}.
We can return an exact feasible point here using an output
representation \(a+b\sqrt d\),
giving scalar \(d\ge0\) and vectors \(a,b\)
in the field generated by the input.
\end{corollary}


\begin{proof}
In the single-constraint case,
test first for feasibility with \(h(x)=0\), by LP
(since we then just require \(F(x)=0\) for~\eqref{eq:lor}).
For \(h(x)>0\) use variables \(s=1/h(x),y=sx\),
writing the conditions as \(s>0\)
together with
\[
 Ay\le sb,\qquad H(y,s)=1,\qquad f(y,s):=\|T(y,s)\|_2^2\le1,
\]
with \(Ax\le b\) the given linear constraints and \(H,T\)
the linear forms or maps obtained by homogenizing
\(h,F\) using \(s\).
Find a point satisfying all but possibly the norm constraint
(or conclude infeasibility) by LP.
If this point does not already give a solution,
minimize \(f\) over the linear conditions with \(s>0\)
relaxed to \(s\ge0\),
by Corollary~\ref{cor:qp}.
A minimum greater than one proves infeasibility. If the minimum
is less than one, interpolate between its value and the value at the
point from LP, choosing a positive weight on the latter so the
interpolated value is one. The corresponding convex combination has
\(s>0\) and \(f\le1\) by convexity. If the minimum is one, we need
a minimizer with \(s>0\). Strict convexity of the squared norm in
the \(T\)-image makes that image the same at every minimizer.
Fix it to the value already found and test for \(s>0\) by LP.


Now for several constraints of the indicated type,
require the right-hand sides nonnegative and square the constraints.
Besides linear constraints the result has the form
\[
 \ell_j(x)\le q(x),\qquad q(x)\le r_k(x)\qquad(\text{respective indices }j,k)
\]
for affine \(\ell_j,r_k\) and one quadratic form \(q\),
after rescaling inequalities (all quadratic parts
are multiples of one).
Impose also all \(\ell_j(x)\le r_k(x)\).

On this linear region, test separately for feasibility with the
lower bounds, and for feasibility with the upper bounds.
For the lower bounds, try each \(j\) as a maximizing index. Impose
\(\ell_j\ge\ell_{j'}\) for all \(j'\) and enforce \(\ell_j\le q\)
using its original Lorentz constraint. The upper bounds are treated
in the same way by enumerating a minimizing index. If either list
is empty, its separate feasibility test consists of linear
constraints alone.
If there are both kinds, with feasibility for each separately,
consider the segment between the two witnesses. At each point, the lower
bounds all hold or the upper bounds all hold (since \(\ell_j\le r_k\)).
These are closed conditions, so somewhere on the segment both hold.
We can find such a point by testing endpoints and enumerating roots
of nonconstant quadratic or linear constraint polynomials on the
segment (an endpoint of the closed feasible portion suffices).
Testing at such a root takes ordinary scalar operations and comparisons
by isolating a radical and squaring (checking signs first).
Only constantly many layers of subroutine calls
or polynomial-size arithmetic formulas are needed,
so we get strong polynomiality as claimed.
\end{proof}

A radical in the generalization can be necessary for a witness.
For example, \(\|(2x,1)\|_2\le3\) and
\(\|(1,1)\|_2\le x\) together force \(x=\sqrt2\).
